% =============================================================================
\section{Introduction}
\label{sec:intro}
% =============================================================================

Three variables shape production LLM inference: the workload, the
routing strategy, and the GPU pool architecture.  These are studied by
largely separate communities---workload characterization, model
routing, and systems/fleet optimization---that seldom interact.  But in
practice, a routing decision tuned for one workload on one pool
topology can be suboptimal for a different combination.

\paragraph{Our research program.}
The vLLM Semantic Router (\vsr{}) project~\cite{vllmsr2026} has been
building the \emph{inference routing stack}.  Across a growing body
of work (summarized in \Cref{tab:ownwork} and
\Cref{sec:foundation}), we have addressed four pillars:
\begin{itemize}[nosep]
  \item \textbf{Routing architecture}: signal-driven semantic routing
        with composable signals~\cite{vllmsr2026}, conflict-free
        policy languages for probabilistic ML
        predicates~\cite{probpol2026}, 98$\times$ router latency
        reduction via Flash Attention and prompt
        compression~\cite{fastrouter2026}, 2D Matryoshka embeddings
        for configurable latency--quality
        trade-offs~\cite{mmberted2025}, reasoning-selective
        routing~\cite{wang2025reason}, category-aware semantic
        caching~\cite{wang2025cache}, feedback-driven online
        routing adaptation~\cite{feedbackdet2026},
        hallucination-aware response validation with conditional
        factcheck classification and token-level
        detection~\cite{halugate2025, factcheck2026}, and
        hierarchical content-safety classification following the
        MLCommons AI Safety Taxonomy with a binary safe/unsafe
        gate~\cite{safetyl1_2026} and a 9-class hazard
        categorizer~\cite{safetyl2_2026} enabling per-category policy
        enforcement and privacy-preserving routing;
  \item \textbf{Pool routing and fleet optimization}: token-budget pool
        routing~\cite{poolrouting2026}, FleetOpt analytical fleet
        provisioning with gateway-layer
        compression~\cite{fleetopt2026}, inference-fleet-sim for
        queueing-grounded capacity
        planning~\cite{fleetsim2026}, and the 1/W law for energy
        efficiency~\cite{onewlaw2026};
  \item \textbf{Multimodal and agent routing}: adaptive VLM routing for
        computer use agents~\cite{avr2026}, outcome-aware tool
        selection~\cite{oats2026}, the Visual Confused Deputy
        security guardrail~\cite{vcd2026}, and the
        gateway-centric personal-assistant stack
        OpenClaw~\cite{openclaw2026} (open-source; multi-channel
        sessions and tools behind a single Gateway control plane);
  \item \textbf{Governance and standards}: the Semantic Inference
        Routing Protocol (SIRP) at IETF~\cite{sirp2025} and
        multi-provider extensions for agentic AI inference
        APIs~\cite{mpext2025}.
\end{itemize}

Each paper solved a specific problem, but the problems are coupled.  
For example, pool routing~\cite{poolrouting2026} creates a cost
cliff whose resolution via gateway-layer compression depends on the
workload's prompt-length CDF archetype~\cite{fleetopt2026}.  The 1/W
law~\cite{onewlaw2026} shows that routing topology is a stronger
energy lever than GPU generation, but only when the pool supports
context-length partitioning.

\begin{figure}[t]
      \centering
      \begin{tikzpicture}[
          >=Stealth,
          box/.style={draw, rounded corners=3pt, minimum width=3.2cm,
                      minimum height=1cm, font=\small\sffamily, align=center,
                      line width=0.8pt},
          dimbox/.style={box, line width=1.2pt, minimum width=4cm,
                         minimum height=1.3cm, font=\sffamily\bfseries},
          arr/.style={->, line width=0.8pt},
          darr/.style={<->, line width=0.8pt, dashed},
      ]
      % Dimension boxes
      \node[dimbox, fill=workloadblue!15, draw=workloadblue]
          (W) at (0, 0) {Workload\\Characterization};
      \node[dimbox, fill=routergreen!15, draw=routergreen]
          (R) at (5.5, 0) {Router\\Layer};
      \node[dimbox, fill=poolorange!15, draw=poolorange]
          (P) at (11, 0) {Pool\\Architecture};
      
      % Sub-items for Workload
      \node[box, fill=workloadblue!5, draw=workloadblue!60,
            minimum width=3.6cm, font=\scriptsize\sffamily]
          (w1) at (0, -2) {Chat vs.\ Agent};
      \node[box, fill=workloadblue!5, draw=workloadblue!60,
            minimum width=3.6cm, font=\scriptsize\sffamily]
          (w2) at (0, -3.2) {Warm vs.\ Cold};
      \node[box, fill=workloadblue!5, draw=workloadblue!60,
            minimum width=3.6cm, font=\scriptsize\sffamily]
          (w3) at (0, -4.4) {Prefill-heavy vs.\ Decode-heavy};
      
      % Sub-items for Router
      \node[box, fill=routergreen!5, draw=routergreen!60,
            minimum width=3.6cm, font=\scriptsize\sffamily]
          (r1) at (5.5, -2) {Static semantic rules};
      \node[box, fill=routergreen!5, draw=routergreen!60,
            minimum width=3.6cm, font=\scriptsize\sffamily]
          (r2) at (5.5, -3.2) {Online bandit / feedback};
      \node[box, fill=routergreen!5, draw=routergreen!60,
            minimum width=3.6cm, font=\scriptsize\sffamily]
          (r3) at (5.5, -4.4) {RL / cascading selection};
      
      % Sub-items for Pool
      \node[box, fill=poolorange!5, draw=poolorange!60,
            minimum width=3.6cm, font=\scriptsize\sffamily]
          (p1) at (11, -2) {Homo vs.\ Hetero GPU};
      \node[box, fill=poolorange!5, draw=poolorange!60,
            minimum width=3.6cm, font=\scriptsize\sffamily]
          (p2) at (11, -3.2) {Disaggregated P/D};
      \node[box, fill=poolorange!5, draw=poolorange!60,
            minimum width=3.6cm, font=\scriptsize\sffamily]
          (p3) at (11, -4.4) {KV-cache topology};
      
      % Arrows from dimension to sub-items
      \foreach \s in {w1,w2,w3} \draw[arr, workloadblue!70] (W) -- (\s);
      \foreach \s in {r1,r2,r3} \draw[arr, routergreen!70] (R) -- (\s);
      \foreach \s in {p1,p2,p3} \draw[arr, poolorange!70] (P) -- (\s);
      
      % Cross-dimensional interactions
      \draw[darr, objectivepurple] (W) -- (R)
          node[midway, above=1pt, font=\scriptsize\itshape, text=objectivepurple]
          {workload shapes routing};
      \draw[darr, objectivepurple] (R) -- (P)
          node[midway, above=1pt, font=\scriptsize\itshape, text=objectivepurple]
          {routing shapes pool need};
      \draw[darr, objectivepurple] (W) to[bend left=25] node[midway, above=2pt,
          font=\footnotesize\itshape, text=objectivepurple] {workload determines pool sizing} (P);
      
      % Safety & Privacy cross-cutting bar
      \definecolor{safetyred}{RGB}{180,40,60}
      \node[draw=safetyred, fill=safetyred!8, rounded corners,
            minimum width=12cm, minimum height=0.8cm,
            font=\small\sffamily\bfseries, text=safetyred]
          (safety) at (5.5, -5.8) {Safety \& Privacy: Content Gate
          \textbullet\ Jailbreak Detection
          \textbullet\ Hazard Classification
          \textbullet\ Privacy Routing};
      
      \draw[arr, safetyred!50] (safety) -- (0, -5.1);
      \draw[arr, safetyred!50] (safety) -- (5.5, -5.1);
      \draw[arr, safetyred!50] (safety) -- (11, -5.1);
      
      % Objective bar
      \node[draw=objectivepurple, fill=objectivepurple!8, rounded corners,
            minimum width=12cm, minimum height=0.8cm,
            font=\small\sffamily\bfseries, text=objectivepurple]
          (obj) at (5.5, -7.2) {Optimization Objectives: Cost
          \textbullet\ Accuracy \textbullet\ Latency \textbullet\ Energy};
      
      \draw[arr, objectivepurple!50] (obj) -- (0, -6.3);
      \draw[arr, objectivepurple!50] (obj) -- (5.5, -6.3);
      \draw[arr, objectivepurple!50] (obj) -- (11, -6.3);
      \end{tikzpicture}
      \caption{The Workload--Router--Pool (WRP) architecture.  Three
      dimensions interact through dashed bidirectional arrows.  Safety \&
      Privacy is a cross-cutting layer: content gating and hazard
      classification shape all three dimensions.  All are further
      constrained by the optimization objectives at the bottom.  Our
      prior work addresses one or more cells in this framework.}
      \label{fig:wrp}
      \end{figure}

      
\paragraph{Our contributions.}
We present the \textbf{Workload--Router--Pool (WRP) architecture}
(\Cref{fig:wrp}), a three-dimensional framework that organizes both
our prior work and the broader research landscape.  The
contributions are:
\begin{enumerate}[nosep]
  \item A \textbf{structural decomposition} into three interacting
        dimensions---Workload, Router, Pool---grounded in evidence from
        our previous work.
  \item A \textbf{mapping} of our prior work onto the
        $3 \times 3$ WRP interaction matrix, showing which cells we
        have addressed and which remain open.
  \item A \textbf{research roadmap} of twenty-one concrete opportunities at
        the intersections, each grounded in our prior measurements.
\end{enumerate}


